% ============================================================
\newpage
\section{Phân tích Hiệu năng và Capacity Planning}
\label{sec:performance}
% ============================================================

\subsection{Tổng quan}

Hệ thống Data Integration Hub vận hành theo hai chiều dữ liệu với
đặc tính tải hoàn toàn khác nhau:

\begin{itemize}
  \item \textbf{Inbound (Sync Job)}: Lấy dữ liệu từ Source API và
        ghi vào PostgreSQL theo lịch định kỳ. Tải có thể dự đoán,
        bottleneck chính là \textbf{RAM} và \textbf{network I/O}.
  \item \textbf{Outbound (Export API)}: Phục vụ các hệ thống tiêu
        thụ truy vấn dữ liệu qua REST API phân trang. Tải không dự
        đoán được, bottleneck chính là \textbf{DB query latency} và
        \textbf{concurrency}.
\end{itemize}

Để kiểm tra hiệu năng có kiểm soát, nhóm xây dựng một
\textbf{Mock Source Server} (Express.js, port 3001) cung cấp các
endpoint mô phỏng nhiều kịch bản dữ liệu khác nhau. Server này
phục vụ đồng thời hai mục đích: kiểm thử cơ chế phòng thủ
\textit{Data Collision} và đo lường hiệu năng pipeline.

% ------------------------------------------------------------
\subsection{Kiểm thử Data Collision — Mock Source Server}
\label{sec:perf-collision-test}
% ------------------------------------------------------------

\subsubsection{Kiến trúc Mock Server}

Mock server được thiết kế với một helper dùng chung
\texttt{buildResponse()} chuẩn hóa mọi response theo đúng API
Contract mà NestJS kỳ vọng:

\begin{lstlisting}[
  caption={Helper \texttt{buildResponse()} của Mock Server},
  label={lst:build-response}]
function buildResponse(tableName, allData, page, limit,
                       overrideCurrentPage = null) {
  const totalRecords = allData.length;
  const totalPages   = Math.ceil(totalRecords / limit) || 1;
  const startIndex   = (page - 1) * limit;
  const payload      = allData.slice(startIndex, startIndex + limit);

  return {
    success:   true,
    timestamp: new Date().toISOString(),
    metadata:  { tableName, totalRecords, totalPages,
                 currentPage: overrideCurrentPage ?? page, limit },
    payload,
  };
}
\end{lstlisting}

\subsubsection{Ma trận Test Case}

Bảng~\ref{tab:collision-testcases} tổng hợp toàn bộ 15 test case
được thiết kế để kiểm tra từng lớp cơ chế phòng thủ. Mỗi test
case trỏ vào một endpoint riêng biệt trên Mock Server và được cấu
hình qua Schema Registry MongoDB.

\begin{longtable}{cp{3cm}p{5cm}p{4.5cm}}
\caption{Ma trận kiểm thử Data Collision Defense}
\label{tab:collision-testcases} \\
\toprule
\textbf{TC} & \textbf{Tên} & \textbf{Kịch bản / Cơ chế kiểm tra} & \textbf{Kỳ vọng} \\
\midrule
\endfirsthead
\multicolumn{4}{r}{\textit{(tiếp theo)}} \\
\toprule
\textbf{TC} & \textbf{Tên} & \textbf{Kịch bản / Cơ chế kiểm tra} & \textbf{Kỳ vọng} \\
\midrule
\endhead
\midrule
\multicolumn{4}{r}{\textit{(xem trang sau)}} \\
\endfoot
\bottomrule
\endlastfoot

TC-01 & Happy Path
  & Dữ liệu chuẩn, đúng contract, không có lỗi
  & Sync 100\% thành công \\[4pt]

TC-02 & Schema Filter
  & Payload có 5 fields lạ không khai báo trong Prisma DMMF
  & Fields bị lọc, upsert thành công \\[4pt]

TC-03 & Dedup
  & id=801 lặp ×3, id=804 lặp ×2 trong 1 batch
  & Giữ record cuối, log \textit{``2 duplicates removed''} \\[4pt]

TC-04 & Type Norm (Int/Bool)
  & \texttt{id="701"} (string), \texttt{active="true"} (string)
  & Normalize thành \texttt{Int} và \texttt{Boolean} đúng kiểu \\[4pt]

TC-04b & Type Norm (DateTime)
  & ISO string, date-only string, chuỗi sai định dạng
  & ISO → Date; chuỗi sai → \texttt{null} (không crash) \\[4pt]

TC-05 & Contract: thiếu \texttt{success}
  & Response dùng key \texttt{data} thay vì \texttt{payload}, thiếu \texttt{success}
  & Throw \textit{[API CONTRACT VIOLATION]}, skip bảng \\[4pt]

TC-06 & Contract: \texttt{success: false}
  & Hệ thống nguồn trả lỗi nội bộ
  & Throw \textit{[API CONTRACT VIOLATION]}, skip bảng \\[4pt]

TC-07 & Contract: payload object
  & \texttt{payload} là \texttt{\{\}} thay vì \texttt{[]}
  & Throw \textit{[API CONTRACT VIOLATION]}, skip bảng \\[4pt]

TC-08 & Contract: thiếu \texttt{payload}
  & Response chỉ có \texttt{metadata}, không có \texttt{payload}
  & Throw \textit{[API CONTRACT VIOLATION]}, skip bảng \\[4pt]

TC-09 & Page Mismatch
  & Server luôn trả \texttt{currentPage: 1} dù request page 2
  & \textit{[PAGE MISMATCH]} → dừng loop ở trang 2 \\[4pt]

TC-10 & Empty Payload
  & \texttt{totalPages: 3} nhưng page 2 trả \texttt{payload: []}
  & \textit{[EMPTY PAYLOAD]} → dừng loop sớm \\[4pt]

TC-11 & Dead Letter: thiếu PK
  & 3/6 records có \texttt{id: null} hoặc không có field \texttt{id}
  & 3 records lỗi vào Dead Letter Log; 3 còn lại sync thành công \\[4pt]

TC-12 & Dead Letter: DB Constraint
  & \texttt{ma = "THIS\_STRING\_} \texttt{EXCEEDS\_VARCHAR5"}
  & Record vi phạm vào Dead Letter Log; batch không bị abort \\[4pt]

TC-13 & Schema Status Gate
  & Set \texttt{status="changed"} trong Schema Registry
  & Bảng bị skip, không fetch API \\[4pt]

TC-14 & Bad Endpoint
  & Set \texttt{dataFromApi} sai đường dẫn trong Schema Registry
  & HTTP 404, retry 3 lần, fail fast \\[4pt]

TC-15 & Idempotent Upsert
  & Chạy sync bảng \texttt{dm\_gioi\_tinh} 2 lần liên tiếp
  & Row count PostgreSQL không tăng ở lần 2 \\

\end{longtable}

\subsubsection{Ví dụ kết quả kiểm thử — TC-03 Deduplication}

\begin{lstlisting}[
  caption={Payload TC-03: id=801 lặp 3 lần trong cùng 1 batch},
  label={lst:tc03-payload}]
const dupData = [
  { id: 801, ma: "D1A", ten: "Dup801 v1 overwrite", active: true  },
  { id: 802, ma: "N01", ten: "Binh thuong #1",       active: true  },
  { id: 801, ma: "D1B", ten: "Dup801 v2 overwrite", active: true  },
  { id: 803, ma: "N02", ten: "Binh thuong #2",       active: false },
  { id: 801, ma: "D1C", ten: "Dup801 v3 KEPT",      active: false }, // <- keep
  { id: 804, ma: "D2A", ten: "Dup804 v1 overwrite", active: true  },
  { id: 804, ma: "D2B", ten: "Dup804 v2 KEPT",      active: false }, // <- keep
];
\end{lstlisting}

Kết quả thực tế từ NestJS log:

\begin{verbatim}
WARN [SyncEngineService] [DEDUP] 2 duplicate records removed (pk: id)
LOG  [SyncEngineService] Sync table dm_gioi_tinh finished.
     Success: 5, Failed: 0
\end{verbatim}

7 records đầu vào → sau dedup còn 5 → upsert 5 thành công. Record
\texttt{id=801} trong PostgreSQL có giá trị \texttt{ma="D1C"} (bản
cuối cùng).

\subsubsection{Ví dụ kết quả kiểm thử — TC-11 Dead Letter}

\begin{lstlisting}[caption={Payload TC-11: mix records hợp lệ và thiếu PK},
  label={lst:tc11-payload}]
const mixedBatch = [
  { id: 401, ma: "OK_1", ten: "Record hop le #1", active: true  },
  { id: null, ma: "ER_1", ten: "PK null",         active: true  }, // <- dead letter
  { id: 402, ma: "OK_2", ten: "Record hop le #2", active: false },
  { /* thieu id */  ma: "ER_2", ten: "No id",     active: true  }, // <- dead letter
  { id: 403, ma: "OK_3", ten: "Record hop le #3", active: true  },
  { id: null, ma: "ER_3", ten: "PK null 2",       active: false }, // <- dead letter
];
\end{lstlisting}

Kết quả thực tế:

\begin{verbatim}
ERROR [SyncEngineService] [FIRST ERROR SAMPLE]
      pk=null | Missing primary key (id) in record
WARN  [SyncEngineService] [DEDUP] 3 records with null/undefined PK detected
LOG   [SyncEngineService] Sync table dm_gioi_tinh finished.
      Success: 3, Failed: 3
\end{verbatim}

3 records lỗi được lưu vào \texttt{errors[]} trong MongoDB EventLog
với chi tiết \texttt{\{pk: null, error: "Missing primary key"\}}.
Batch không bị abort — 3 records hợp lệ vẫn được sync thành công.

% ------------------------------------------------------------
\subsection{Kiểm thử Hiệu năng — PERF Test Cases}
\label{sec:perf-ram}
% ------------------------------------------------------------

\subsubsection{Thiết kế MemoryProfiler}

Để đo RAM thực tế, nhóm implement \texttt{MemoryProfiler} tích hợp
trực tiếp vào \texttt{syncOneTable()}. Profiler lấy mẫu
\texttt{process.memoryUsage().rss} định kỳ 300ms trong suốt vòng
đời của pipeline:

\begin{lstlisting}[caption={\texttt{MemoryProfiler} đo peak RSS trong sync pipeline},
  label={lst:memory-profiler}]
export class MemoryProfiler {
  private samples: number[] = [];
  private interval: NodeJS.Timeout;

  start(intervalMs = 500) {
    this.samples = [];
    this.interval = setInterval(() => {
      this.samples.push(process.memoryUsage().rss);
    }, intervalMs);
  }

  stop() {
    clearInterval(this.interval);
    const mb = (b: number) => +(b / 1024 / 1024).toFixed(1);
    return {
      peakMB:   mb(Math.max(...this.samples)),
      baseMB:   mb(this.samples[0]),
      deltaRSS: mb(Math.max(...this.samples) - this.samples[0]),
      samples:  this.samples.length,
    };
  }
}
\end{lstlisting}

Profiler được khởi động ngay đầu \texttt{syncOneTable()} và dừng
trong block \texttt{finally} để đảm bảo luôn có kết quả kể cả khi
pipeline thất bại.

\subsubsection{PERF-01 — RAM Capacity Planning}

Mock endpoint \texttt{/perf/ram-test} sinh dữ liệu tổng hợp với
kích thước record kiểm soát được qua tham số \texttt{records} và
\texttt{recordSize}. Padding được nhồi vào field \texttt{hinhThePath}
(\texttt{VarChar(500)}) của bảng \texttt{nguoi\_hoc} để không vi
phạm Schema Contract:

\begin{lstlisting}[caption={PERF-01: Sinh dữ liệu tổng hợp với kích thước kiểm soát},
  label={lst:perf01-mock}]
// GET /perf/ram-test?records=10000&recordSize=500&page=1&limit=5000
const paddingLen = Math.min(Math.max(0, recordSize - 120), 480);
const pathPad    = "p".repeat(paddingLen);

const allData = Array.from({ length: totalRecords }, (_, i) => ({
  id:          10000 + i,
  ho:          "Nguyen",
  ten:         `SinhVien${i}`,
  gioiTinh:    i % 2 === 0 ? "M" : "F",
  emailCaNhan: `sv${i}@hcmut.edu.vn`,
  hinhThePath: `/photos/sv${i}.jpg${pathPad}`,  // <- padding
  updatedAt:   new Date(...).toISOString(),
}));
\end{lstlisting}

\textbf{Kết quả thực nghiệm} — 10.000 records, recordSize = 500 bytes,
batch size = 5.000:

\begin{verbatim}
[MEMORY] nguoi_hoc: peak=775.9MB, base=224.4MB, delta=551.6MB
[DONE]   nguoi_hoc: synced=10000, skipped(unchanged)=0.
\end{verbatim}

\begin{table}[H]
\centering
\caption{Phân tích nguồn tiêu thụ RAM trong pipeline (PERF-01)}
\label{tab:ram-breakdown}
\begin{tabularx}{\textwidth}{lrX}
\toprule
\textbf{Thành phần} & \textbf{Ước lượng} & \textbf{Ghi chú} \\
\midrule
Pre-sync Backup
  & $\approx 340$ MB
  & Serialize 68.411 records từ PostgreSQL thành JSON \\
Sync 2 trang (5k records/trang)
  & $\approx 45$ MB
  & V8 parse overhead $\times 3$, giải phóng sau mỗi trang \\
Orphan Detection
  & $\approx 25$ MB
  & \texttt{findMany(\{select: \{id\}\})} — 78k IDs \\
HTTP + EventEmitter overhead
  & $\approx 30$ MB
  & Axios buffer, event queue \\
\midrule
\textbf{Tổng delta đo được} & \textbf{551,6 MB} & RSS thực tế \\
\bottomrule
\end{tabularx}
\end{table}

\textbf{Nhận xét quan trọng}: Delta 551 MB không phản ánh chi phí
của việc sync 10.000 records mà phản ánh toàn bộ pipeline, trong đó
\textbf{backup chiếm phần lớn} do phải serialize 68.411 records có
sẵn trong DB thành JSON buffer trước khi upload MinIO. Chi phí
thuần của sync (2 trang × 5.000 records) chỉ khoảng 45 MB.

\subsubsection{PERF-02 — Incremental Sync với Server-side Filter}

Mock endpoint \texttt{/perf/incremental-test} tạo sẵn 1.000 records
với \texttt{updatedAt} trải đều trong 30 ngày qua. Khi nhận tham số
\texttt{updatedAfter}, server lọc và chỉ trả về records thỏa điều kiện:

\begin{lstlisting}[caption={PERF-02: Server-side filter theo \texttt{updatedAfter}},
  label={lst:perf02-mock}]
// GET /perf/incremental-test?updatedAfter=2026-04-22T00:00:00Z
const filtered = updatedAfter
  ? allData.filter(r => new Date(r.updatedAt) > updatedAfter)
  : allData; // lan dau sync -> full load

// Log mock server:
// total=1000, filtered=40 (4.0% tra ve) | page=1/1
\end{lstlisting}

\textbf{Kết quả thực nghiệm} — \texttt{updatedAfter=2026-04-22T00:00:00Z},
\texttt{lastSyncTime} lấy từ Schema Registry MongoDB:

\begin{verbatim}
[STRATEGY: INCREMENTAL] nguoi_hoc
  — server-side filter updatedAfter=2026-04-22T00:00:00.000Z
Fetching page 1/1 -> .../perf/incremental-test?...
    &updatedAfter=2026-04-22T00%3A00%3A00.000Z
Table nguoi_hoc has 40 records across 1 pages.
[INCREMENTAL] Page 1: 40 records thay đổi từ source — upsert.
Sync table nguoi_hoc finished. Success: 40, Failed: 0
[MEMORY] nguoi_hoc: peak=831.2MB, base=212MB, delta=619.3MB
\end{verbatim}

\begin{table}[H]
\centering
\caption{So sánh \texttt{upsert} vs \texttt{incremental} — kết quả thực nghiệm}
\label{tab:perf-comparison}
\begin{tabular}{lcc}
\toprule
\textbf{Chỉ số} & \textbf{PERF-01 (upsert)} & \textbf{PERF-02 (incremental)} \\
\midrule
Records fetch qua network   & 10.000            & \textbf{40} \\
Records ghi DB (upsert)     & 10.000            & \textbf{40} \\
Số trang                    & 2                 & \textbf{1} \\
Thời gian sync thuần        & $\approx 9$ s     & $\approx 1$ s \\
DB writes tiết kiệm         & —                 & \textbf{99,6\%} \\
Network tiết kiệm           & —                 & \textbf{99,6\%} \\
\midrule
Delta RAM                   & 551,6 MB          & 619,3 MB \\
\multicolumn{3}{l}{\footnotesize\textit{Delta RAM của incremental cao hơn vì backup lúc này
phải serialize nhiều records hơn}} \\
\multicolumn{3}{l}{\footnotesize\textit{(78.411 records thay vì 68.411),
không phải do phần sync tốn nhiều hơn.}} \\
\bottomrule
\end{tabular}
\end{table}

Kết quả xác nhận: với \texttt{lastSyncTime = 2026-04-22T00:00:00Z},
chỉ \textbf{4\%} dataset (40/1.000 records) được truyền và ghi,
tiết kiệm 96\% chi phí so với full \texttt{upsert}.

% ------------------------------------------------------------
\subsection{Phân tích RAM cho Sync Job}
\label{sec:perf-ram-analysis}
% ------------------------------------------------------------

\subsubsection{Vòng đời dữ liệu trong RAM}

Kết quả từ PERF-01 và PERF-02 cho thấy pipeline tiêu thụ RAM
qua bốn giai đoạn nối tiếp:

\begin{figure}[H]
\centering
\begin{tikzpicture}[
  stage/.style={rectangle, draw, rounded corners, minimum width=2.5cm,
                minimum height=1.4cm, text centered, align=center, font=\small},
  arr/.style={-{Stealth}, thick},
  lbl/.style={font=\footnotesize, text=gray},
]
\node[stage, fill=orange!20] (bk)    {Pre-sync\\Backup\\(MinIO)};
\node[stage, fill=blue!10,   right=0.8cm of bk]   (net)   {HTTP\\Response\\Buffer};
\node[stage, fill=green!10,  right=0.8cm of net]  (parse) {JSON Parse\\+ Normalize\\(V8)};
\node[stage, fill=yellow!10, right=0.8cm of parse](db)    {Prisma\\upsert\\per-row};
\node[stage, fill=red!10,    right=0.8cm of db]   (orp)   {Orphan\\Detection\\(IDs only)};

\draw[arr] (bk) -- (net);
\draw[arr] (net) -- (parse);
\draw[arr] (parse) -- (db);
\draw[arr] (db) -- (orp);

% \node[lbl, below=0.3cm of bk]    {Bottleneck chính};
% \node[lbl, below=0.3cm of net]   {per page};
% \node[lbl, below=0.3cm of parse] {tạm thời};
% \node[lbl, below=0.3cm of db]    {nhỏ};
% \node[lbl, below=0.3cm of orp]   {ID only};
\end{tikzpicture}
\caption{Các giai đoạn tiêu thụ RAM trong một lần \texttt{syncOneTable()}.
Chi tiết ước lượng từng thành phần xem Bảng~\ref{tab:ram-breakdown}.}
\label{fig:ram-lifecycle}
\end{figure}

\subsubsection{Benchmark matrix RAM theo số lượng records}

Để đánh giá quy luật tăng trưởng RAM, nhóm chạy PERF-01 với các
cấu hình khác nhau về số lượng records và kích thước record. Kết quả
thực đo (đánh dấu \textsuperscript{*}) và ước lượng ngoại suy từ
mô hình phân tích được tổng hợp trong Bảng~\ref{tab:ram-matrix}.

\begin{table}[H]
\centering
\caption{Benchmark matrix: delta RAM theo số records và kích thước
(thực đo\textsuperscript{*} và ước lượng ngoại suy)}
\label{tab:ram-matrix}
\begin{tabular}{rrrrrr}
\toprule
\textbf{Records} & \textbf{Size/rec} & \textbf{Batch} &
\textbf{Backup RAM} & \textbf{Sync RAM} & \textbf{Delta tổng} \\
\midrule
1.000   & 200 B & 1.000 & $\approx 170$ MB & $\approx 3$ MB  & $\approx 173$ MB \\
5.000   & 200 B & 5.000 & $\approx 170$ MB & $\approx 7$ MB  & $\approx 177$ MB \\
10.000  & 500 B & 5.000 & $\approx 340$ MB & $\approx 45$ MB & \textbf{551,6 MB\textsuperscript{*}} \\
10.000  & 200 B & 5.000 & $\approx 340$ MB & $\approx 20$ MB & $\approx 360$ MB \\
50.000  & 200 B & 5.000 & $\approx 850$ MB & $\approx 20$ MB & $\approx 870$ MB \\
100.000 & 200 B & 5.000 & $\approx 1,7$ GB & $\approx 20$ MB & $\approx 1,72$ GB \\
\bottomrule
\multicolumn{6}{l}{\footnotesize\textsuperscript{*}~Số liệu thực đo;
  các dòng còn lại là ước lượng ngoại suy theo mô hình tuyến tính.} \\
\multicolumn{6}{l}{\footnotesize Backup RAM $\propto$ records trong DB
  tại thời điểm sync; Sync RAM $\propto$ recordSize $\times$ batchSize.}
\end{tabular}
\end{table}

Nhận xét quan trọng từ ma trận trên: \textbf{Backup RAM tăng tuyến
tính theo số records trong DB} và là thành phần chi phối, trong khi
Sync RAM gần như hằng số theo batch (không phụ thuộc tổng số records
cần sync mà chỉ phụ thuộc \texttt{recordSize × batchSize}).

\begin{table}[H]
\centering
\caption{Ước lượng peak RAM theo cấu hình bảng (ngoại suy từ điểm thực đo)}
\label{tab:ram-estimate}
\begin{tabular}{lrrrr}
\toprule
\textbf{Bảng} & \textbf{DB records} & \textbf{Batch size} &
\textbf{Backup RAM} & \textbf{Sync RAM/batch} \\
\midrule
\texttt{dm\_*} (danh mục)  & $<1.000$   & 1.000  & $\approx 5$ MB   & $\approx 2$ MB \\
\texttt{nguoi\_hoc} (thực) & 68.411     & 5.000  & $\approx 170$ MB & $\approx 10$ MB \\
Bảng lớn giả định         & 500.000    & 5.000  & $\approx 1,2$ GB & $\approx 10$ MB \\
\bottomrule
\multicolumn{5}{l}{\footnotesize Ước lượng ngoại suy; chưa thực đo trực tiếp các cấu hình này.}
\end{tabular}
\end{table}

\textbf{Kết luận}: Với bảng lớn (>100k records), \textbf{backup là
bottleneck RAM}, không phải sync. Hướng cải thiện là stream upload
MinIO theo từng chunk thay vì serialize toàn bộ vào buffer trước.

% ------------------------------------------------------------
\subsection{Phân tích Throughput cho Export API (Outbound)}
\label{sec:perf-export}
% ------------------------------------------------------------

\subsubsection{Kiến trúc Export API hiện tại}

Hệ thống cung cấp dữ liệu ra ngoài qua \texttt{MasterDataController}
với endpoint phân trang:

\begin{verbatim}
GET /api/master-data/:table?page=1&limit=10&search=...
\end{verbatim}

Luồng xử lý mỗi request: \textit{HTTP Client} → \textit{JWT +
Permission Guard} → \textit{MasterDataService} →
\textit{PostgreSQL (Prisma)} → \textit{JSON response}.

\subsubsection{Bottleneck: OFFSET Pagination}

Triển khai hiện tại trong \texttt{MasterDataService.findAll()} sử dụng
\textbf{OFFSET pagination} thông qua Prisma \texttt{skip}:

\begin{lstlisting}[caption={OFFSET pagination trong \texttt{MasterDataService.findAll()}},
  label={lst:offset-pagination}]
const skip = (Number(page) - 1) * Number(limit);
const [total, data] = await Promise.all([
  model.count({ where: whereCondition }),
  model.findMany({ where: whereCondition, skip,
                   take: Number(limit), orderBy: { ma: 'asc' } }),
]);
\end{lstlisting}

OFFSET tương đương SQL: PostgreSQL phải đọc và bỏ qua
\texttt{(page-1)×limit} rows trước khi lấy \texttt{limit} rows
cần thiết — latency tăng tuyến tính theo độ sâu trang:

\begin{table}[H]
\centering
\caption{Ước lượng latency theo độ sâu trang — OFFSET vs Keyset
(số liệu lý thuyết, chưa thực đo trực tiếp)}
\label{tab:pagination-latency}
\begin{tabular}{lcc}
\toprule
\textbf{Page} & \textbf{OFFSET (hiện tại)} & \textbf{Keyset (đề xuất)} \\
\midrule
1    &  $\approx$45 ms   & $\approx$45 ms \\
50   & $\approx$120 ms   & $\approx$47 ms \\
100  & $\approx$380 ms   & $\approx$47 ms \\
500  & $\approx$2.100 ms & $\approx$48 ms \\
\bottomrule
\multicolumn{3}{l}{\footnotesize Ước lượng lý thuyết dựa trên độ phức
tạp $O(n)$ của OFFSET và $O(1)$ của Keyset pagination.} \\
\multicolumn{3}{l}{\footnotesize Với bảng \texttt{dm\_*} ($<$1.000 records),
OFFSET không tạo vấn đề thực tế ở quy mô hiện tại.}
\end{tabular}
\end{table}

Với các bảng danh mục nhỏ (\texttt{dm\_*}, $<$1.000 records) hiện tại,
OFFSET không tạo vấn đề. Vấn đề xuất hiện khi bảng \texttt{nguoi\_hoc}
(68k+ records) được query đến các trang sâu.

\subsubsection{Ước lượng Throughput theo Little's Law}

\[
  \text{RPS}_{\max} = \frac{N_{\text{pool}}}{L_{\text{avg (s)}}}
\]

\begin{table}[H]
\centering
\caption{Ước lượng RPS tối đa theo Little's Law
($\text{RPS} = N_\text{pool} / L_\text{avg}$) — số liệu lý thuyết}
\label{tab:rps-estimate}
\begin{tabular}{cccc}
\toprule
\textbf{Pool size} & \textbf{Avg latency} & \textbf{RPS max (lý thuyết)} & \textbf{req/phút} \\
\midrule
10 & 50 ms  & 200 & 12.000 \\
10 & 200 ms & 50  & 3.000  \\
10 & 380 ms & 26  & 1.560  \\
20 & 50 ms  & 400 & 24.000 \\
\bottomrule
\multicolumn{4}{l}{\footnotesize Ước lượng lý thuyết theo Little's Law;
chưa thực đo bằng load testing tool (k6/JMeter).} \\
\multicolumn{4}{l}{\footnotesize Latency thực tế phụ thuộc cấu hình
PostgreSQL, index, và tải đồng thời.}
\end{tabular}
\end{table}

% ------------------------------------------------------------
\subsection{Đánh giá overhead cơ chế bảo mật}
\label{sec:perf-security}
% ------------------------------------------------------------

Ngoài hiệu năng pipeline ETL, nhóm tiến hành đánh giá sơ bộ overhead
của các lớp bảo mật thông qua quan sát thời gian xử lý request trong
NestJS log.

\subsubsection{JWT Guard overhead}

Mỗi request đến API cần đi qua \texttt{JwtAuthGuard} để xác thực
token. Quá trình này bao gồm: đọc header \texttt{Authorization},
decode và verify chữ ký JWT (HS256), sau đó kiểm tra \texttt{exp}.
Từ log NestJS với timestamp, overhead trung bình của guard đo được
khoảng \textbf{1--3 ms} per request trên môi trường phát triển
(Node.js single instance, không cache).

\subsubsection{Permission Guard và RBAC overhead}

\texttt{PermissionGuard} tra cứu bảng phân quyền của user trong
PostgreSQL sau khi JWT pass. Đây là một DB query bổ sung:

\begin{verbatim}
SELECT roleSettings FROM users WHERE id = $userId  -- <1 ms (PK lookup)
\end{verbatim}

Tổng overhead bảo mật (JWT verify + RBAC lookup) ước lượng khoảng
\textbf{2--5 ms} per request, chiếm tỷ lệ nhỏ so với latency
query dữ liệu chính (thường 40--100 ms cho bảng trung bình).

\begin{table}[H]
\centering
\caption{Đánh giá overhead các lớp bảo mật (quan sát từ NestJS log)}
\label{tab:security-overhead}
\begin{tabular}{lll}
\toprule
\textbf{Cơ chế} & \textbf{Overhead đo được} & \textbf{Phương pháp} \\
\midrule
JWT decode + verify (HS256) & 1--3 ms/req & NestJS log timestamp \\
RBAC DB lookup (PK)         & $<$1 ms/req & Prisma query log \\
Kong Gateway rate-limit     & $<$1 ms/req & Kong latency header \\
\midrule
\textbf{Tổng overhead bảo mật} & \textbf{2--5 ms/req} & \\
\midrule
Latency DB query (trang 1)  & 40--100 ms/req & NestJS log \\
\textbf{Tỷ lệ overhead}     & \textbf{$<$10\%} & So với tổng latency \\
\bottomrule
\multicolumn{3}{l}{\footnotesize Đo trên môi trường phát triển (localhost);
chưa thực hiện load test chính thức.}
\end{tabular}
\end{table}

\subsubsection{Kiểm thử cơ chế chặn truy cập trái phép}

Nhóm thực hiện kiểm thử thủ công các trường hợp bảo mật sau:

\begin{table}[H]
\centering
\caption{Kết quả kiểm thử thủ công cơ chế bảo mật}
\label{tab:security-test}
\begin{tabular}{p{5cm}cc}
\toprule
\textbf{Kịch bản} & \textbf{HTTP Status} & \textbf{Kết quả} \\
\midrule
Không có token                    & 401 & Pass \\
Token hết hạn (\texttt{exp} quá)  & 401 & Pass \\
Token sai chữ ký                  & 401 & Pass \\
Role \texttt{reader} gọi sync API & 403 & Pass \\
Role \texttt{writer} gọi admin API& 403 & Pass \\
Vượt rate limit Kong (>100 req/phút) & 429 & Pass \\
SQL injection qua search param   & 400 / dữ liệu an toàn & Pass \\
\bottomrule
\end{tabular}
\end{table}

Tất cả 7 kịch bản kiểm thử đều cho kết quả đúng kỳ vọng. Hạn chế
của đánh giá này là chưa có công cụ load test tự động (k6, OWASP ZAP)
để kiểm tra ở tải cao đồng thời — đây là hướng mở rộng trong tương lai.

% ------------------------------------------------------------
\subsection{Tóm tắt và Khuyến nghị Server}
\label{sec:perf-summary}
% ------------------------------------------------------------

\begin{table}[H]
\centering
\caption{Khuyến nghị cấu hình server cho Data Integration Hub}
\label{tab:server-recommendation}
\begin{tabular}{lll}
\toprule
\textbf{Thành phần} & \textbf{Tối thiểu} & \textbf{Khuyến nghị} \\
\midrule
RAM server              & 1 GB  & 2 GB \\
CPU                     & 1 core & 2 cores \\
DB connection pool      & 10    & 20 (nếu export API tải cao) \\
\texttt{overwrite} bảng & Chỉ dùng cho $<$10k records &
  Không dùng cho bảng lớn \\
Sync job đồng thời      & 1 (thiết kế tuần tự) & 1 \\
\bottomrule
\end{tabular}
\end{table}

Hai điểm cần cải thiện nếu hệ thống scale lên bảng hàng trăm nghìn
records:
\begin{enumerate}
  \item \textbf{Backup streaming}: Thay vì serialize toàn bộ bảng
        vào RAM rồi upload, dùng Prisma cursor + pipe stream trực
        tiếp lên MinIO — giảm peak RAM backup từ hàng trăm MB
        xuống vài chục MB.
  \item \textbf{Keyset pagination} cho Export API: Thay \texttt{skip}
        bằng \texttt{WHERE id > lastSeenId} để latency flat tại
        mọi độ sâu trang.
\end{enumerate}
